\section{Розділ~\ref{sec:muscles}. Вихід на м'язи}

Будову виходу --- рухові одиниці, запас надійності синапсу, увімкнення за розміром, пару м'язів, петлю розтягу й генератори ритму --- виміряно прямо; умова стійкості петлі із затримкою --- теорема; внутрішня модель тіла --- добре обґрунтована гіпотеза; числа рисунків --- розрахунок за моделлю книги.

{\footnotesize
\setlength{\tabcolsep}{3pt}\renewcommand{\arraystretch}{1.15}
\begin{longtable}{>{\raggedright}p{0.5cm}>{\raggedright}p{4.0cm}>{\raggedright}p{5.6cm}>{\raggedright}p{2.3cm}>{\raggedright\arraybackslash}p{1.7cm}}
\toprule
№ & Твердження & Дослід і що виміряно & Джерело & Статус \\
\midrule\endfirsthead
\toprule
№ & Твердження & Дослід і що виміряно & Джерело & Статус \\
\midrule\endhead
1 & Рухова одиниця: один \nt{ACET}-нейрон керує групою волокон, від кількох сотень до пари тисяч; у м'язах тонкого підлаштування одиниці дрібніші & М'язи людини, підрахунок рухових аксонів і м'язових волокон: у середньому близько $340$ волокон на нейрон у міжкістковому м'язі кисті, близько $410$ у плечопроменевому м'язі, близько $1930$ у медіальному литковому. Точного числа для найдрібніших одиниць розділ не наводить. & \cite{Feinstein1955mu} & виміряно \\
2 & Синапс на м'язі вивільняє в кілька разів більше медіатора, ніж потрібно для спрацювання волокна & Огляд вимірювань на нервово-м'язових синапсах: запас надійності (відношення вивільненого медіатора до порогового) у дорослих ссавців зазвичай $3$--$5$; у людини запас тримається більше на складках мембрани одержувача, ніж на величині вивільнення. & \cite{WoodSlater2001} & виміряно \\
3 & Посмикування~\eqref{eq:twitch} наростає за час $T_c$ порядку $50\ms$ & Поодинокі рухові одиниці кисті людини за довільного зусилля, силу усереднено за імпульсами одиниці: час наростання $30$--$100\ms$, у $80\,\%$ одиниць менше за $70\ms$. Функція $(t/T_c)e^{1-t/T_c}$ --- апроксимація з моделі пулу одиниць. & \cite{MilnerBrown1973rec, Fuglevand1993} & виміряно \\
4 & Сума посмикувань~\eqref{eq:forcesum}: за $5$, $15$ і $40$ імп/с сила $0.2$--$1.1F_1$, $1.8$--$2.2F_1$ і близько $5.4F_1$ (рис.~\ref{fig:tetanus}) & Розрахунок за \texttt{models/\allowbreak muscle\_\allowbreak models.py}, $T_c=50\ms$: $0.21$--$1.10$, $1.76$--$2.19$ і $5.28$--$5.49\,F_1$ на останніх $0.3\,\text{с}$. Лінійне додавання --- спрощення: насичення справжнього м'яза в моделі немає. & \texttt{models/\allowbreak muscle\_\allowbreak models.py} & модель \\
5 & Одиниці вмикаються за розміром; найслабші в сотню разів слабші за найсильніші & Вентральні корінці кішки: за наростаючого рефлекторного входу першими розряджаються нейрони з дрібними імпульсами в корінці, тобто малі. Міжкістковий м'яз кисті людини: сила посмикування одиниць від $0.1$ до $10$~г і зростає майже лінійно із силою, за якої одиниця вмикається. & \cite{Henneman1957, MilnerBrown1973rec} & виміряно \\
6 & Крок сили пропорційний силі, $\Delta F/F\approx q-1$~\eqref{eq:sizeprinciple}: рухома кома & Виведено в розділі з геометричної прогресії сил. Узгоджується з дослідом: за великих зусиль на той самий приріст сили вмикається набагато менше нових одиниць. & виведення в розділі; \cite{MilnerBrown1973rec} & теорія \\
7 & Розкид сили зростає пропорційно самій силі & Довільне ізометричне зусилля м'яза великого пальця: стандартне відхилення сили лінійно зростає із середньою силою; за того самого зусилля, викликаного електричною стимуляцією, розкид сталий. Модель пулу: лінійне зростання дає ввімкнення одиниць за порядком сили. & \cite{Jones2002sdn} & виміряно \\
8 & Плавність рухів --- наслідок шуму, що залежить від сигналу & Модель: траєкторія, що мінімізує розкид кінцевої точки за такого шуму, відтворює виміряні траєкторії ока й руки. Прямого досліду, який відрізняв би цю причину від інших, немає. & \cite{HarrisWolpert1998} & гіпотеза \\
9 & Різниця сил пари м'язів задає момент, сума --- жорсткість~\eqref{eq:antagonists} & Теорія керування імпедансом через спільне напруження. Рука людини: жорсткість кожного суглоба, виміряна малими поштовхами, приблизно пропорційна моменту в ньому; різні співвідношення спільного напруження змінюють форму й напрямок еліпса жорсткості кисті. & \cite{Hogan1984, GomiOsu1998stiff} & виміряно \\
10 & Датчик розтягу збуджує свій \nt{ACET}-нейрон і через гальмівного посередника вимикає антагоніста (рис.~\ref{fig:antagonists}) & Спинний мозок кішки: поодинокий посередник, що збуджується волокнами датчиків розтягу, дає моносинаптичні гальмівні потенціали в \nt{ACET}-нейронах м'яза-антагоніста, синаптична затримка $0.28$--$0.42\ms$. Медіатор посередника (гліцин) у цьому досліді не визначали. & \cite{Jankowska1972Ia} & виміряно \\
11 & Затримка петлі: спинномозкової $25$--$30\ms$, через кору в півтора--три рази більша, через око $100$--$200\ms$ & Рука людини, поштовх по суглобу: коротка відповідь м'язів через $20$--$45\ms$, довга $45$--$105\ms$, довільна $120$--$180\ms$. Зміщення видимого положення пальця під час руху: поправка через $160\ms$ у середньому. & \cite{Pruszynski2008rapid, SaundersKnill2003vis} & виміряно \\
12 & $\dd x/\dd t=-Gx(t-d)$ стійке при $Gd<\pi/2$~\eqref{eq:delaystable}; на межі частота $1/(4d)$ & Теорема про корені рівняння із запізненням. Перевірка за \texttt{models/\allowbreak muscle\_\allowbreak models.py}, $d=30\ms$: до $3\,\text{с}$ амплітуда $3\cdot10^{-6}$ при $G=40$, $0.13$ при $G=50$, $38$ при $G=55$ за секунду (межа $52$). & \cite{Hayes1950} & теорія \\
13 & Фізіологічне тремтіння $8$--$12$~Гц; петля розтягу --- одне з його джерел & Огляд вимірювань: дрібне тремтіння в діапазоні близько $10$~Гц є у здорових людей; його походження змішане --- центральні коливання, властивості розрядів одиниць, механічний резонанс і резонанс рефлекторної петлі, і за різних умов переважає різне. & \cite{McAuleyMarsden2000} & гіпотеза \\
14 & Мозок передбачає результат команди за внутрішньою моделлю тіла & Людина тримає вантаж щипком і рухає рукою: за інерційного, в'язкого й пружного навантаження сила стискання змінюється одночасно з навантаженням, а не із запізненням петлі, тобто навантаження передбачено. Огляд зводить такі досліди; прямого спостереження самої моделі в нейронах немає. & \cite{FlanaganWing1997grip, WolpertGhahramani2000} & гіпотеза \\
15 & Ритм ходьби, дихання, жування породжують місцеві мережі за сталого входу & Огляд дослідів: ізольована нервова система (спинний мозок, ганглії) видає ритмічний руховий візерунок без ритмічного входу й без зворотного зв'язку від м'язів. & \cite{MarderBucher2001} & виміряно \\
16 & Взаємне гальмування з утомою~\eqref{eq:halfcentre} дає ритм із періодом $0.84\,\text{с}\approx2\tau_a$ (рис.~\ref{fig:halfcentre}) & Теорема: мережа із взаємним гальмуванням і адаптацією без стійкого стану за сталого входу обов'язково коливається. Розрахунок за \texttt{models/\allowbreak muscle\_\allowbreak models.py} ($I=1$, $\beta=2$, $b=2$, $\tau=20\ms$, $\tau_a=0.4\,\text{с}$): період $0.84\,\text{с}$. Що справжні генератори влаштовані саме так, не показано. & \cite{Matsuoka1985osc} & модель \\
\bottomrule
\end{longtable}}
